\documentclass[12pt,a4paper]{ctexart}

\usepackage{geometry}
\usepackage{graphicx}
\usepackage{array}
\usepackage{booktabs}
\usepackage{xeCJK}
\usepackage{tikz}
\usepackage{titlesec}

% 页面设置
\geometry{
    top=2.5cm,
    bottom=2.5cm,
    left=2.8cm,
    right=2.8cm
}

% 定义 section* 格式
\titleformat{\section}
  {\zihao{-3}\heiti\bfseries}
  {}
  {0em}
  {}
\titlespacing*{\section}{0pt}{0.5cm}{0.3cm}  % 左缩进、上间距、下间距


% 取消页码（封面）
\pagestyle{empty}

% 定义等长下划线命令（固定宽度 6.5cm，左对齐）
\newcommand{\infofill}[1]{{\zihao{4}\songti\underline{\makebox[6.5cm][c]{#1}}}}

% 定义内容区块命令
\newcommand{\reportsection}[2]{%
    \noindent{\zihao{-3}\heiti\bfseries #1（#2分）}
    \vspace{0.3cm}
    
    \noindent\rule{\textwidth}{0.4pt}
    \vspace{6cm}
}

% 设置中文字体
\setCJKmainfont{SimSun}  % 宋体，可根据系统改为 Noto Serif CJK SC 等

\begin{document}
% ==================== 封面页 ====================
\begin{titlepage}
    \centering
    \vspace*{1.5cm}
    % 校徽占位（正方形示例图）
    \includegraphics[width=5cm,height=5cm]{images/image-2026-09-07-14-57-04.png}
    
    \vspace{1.8cm}
    
    % 中山大学
    {\zihao{1}\heiti\bfseries 中山大学}
    
    \vspace{0.9cm}
    
    % 微电子科学与技术学院
    {\zihao{2}\heiti 微电子科学与技术学院}
    
    \vspace{0.9cm}
    
    % 课程名称
    {\zihao{2}\heiti 《集成电路元器件特性表征实验》}
    
    \vspace{0.9cm}
    
    % 实验报告
    {\zihao{1}\heiti\bfseries 实验报告}
    
    \vspace{3cm}
    
    % 信息表格
    

    % 信息表格（无任何框线，无横线，无竖线）
    \begin{tabular}{@{}>{\zihao{4}\heiti}r @{\hspace{1.2em}} l@{}}
        实验主题      & \infofill{}        \\[0.8cm]
        姓\quad\quad 名 & \infofill{徐睿琳}   \\[0.8cm]
        学\quad\quad 号 & \infofill{23342107} \\[0.8cm]
        实验日期      & \infofill{\today}   \\[0.8cm]
    \end{tabular}
    
    \vfill
\end{titlepage}

% ==================== 内容页 ====================
\newpage
\pagestyle{plain}
\setcounter{page}{2}

\section*{实验目的（5分）}


\section*{实验原理（10分）}


\section*{实验设备和步骤（5分）}


\section*{实验数据分析（50分）}

\section*{实验结论（20分）}


\section*{实验启示和拓展（10分）}



\end{document}